\documentclass[12pt,a4paper]{article}

\usepackage[utf8]{inputenc}
\usepackage{amsmath,amssymb}
\usepackage{graphicx}
\usepackage{hyperref}
\usepackage{booktabs}
\usepackage{geometry}
\usepackage{cite}

\geometry{margin=1in}

% --- Auto-injected Metrics ---
% We will have Python generate a metrics.tex file containing commands like:
% \newcommand{\OptimalAlpha}{0.42}
\IfFileExists{tables/metrics.tex}{\newcommand{\OptimalAlpha}{0.050}
\newcommand{\AlphaRMSE}{0.5850}
}{
    \newcommand{\OptimalAlpha}{TBD}
    \newcommand{\BacktestSpearman}{TBD}
}

\title{Network Contagion of State Instability: A Multiplex Spectral Approach}
\author{Author Name \\ \textit{Institution/Affiliation}}
\date{\today}

\begin{document}

\maketitle

\begin{abstract}
Predicting systemic instability requires moving beyond isolated national indicators to understand cross-border contagion. In this paper, we propose a multiplex network model of state instability propagation based on the Friedkin-Johnsen opinion dynamics framework. By coupling 125 nations through five channels—bilateral trade, financial exposure, geographic proximity, political alignment, and human migration—we demonstrate that network effects significantly amplify latent structural vulnerabilities. The model's propagation coefficient $\alpha$ is empirically calibrated via leave-one-out cross-validation to $\OptimalAlpha$. Backtesting on the 2022 geopolitical shock demonstrates that the multiplex network captures cross-country variance in instability transmission significantly better than geography-only or trade-only baselines ($\rho = \BacktestSpearman$).
\end{abstract}

\section{Introduction}
\input{sections/01_intro.tex}

\section{Methodology}
\input{sections/02_method.tex}

\section{Results}
\input{sections/03_results.tex}

\section{Discussion and Limitations}
\input{sections/04_discussion.tex}

\bibliographystyle{plain}
\bibliography{../references}

\end{document}
